% Terminal note of the homily: everything the speech itself must not carry ---
% the audience, the spoken word count and the pace it implies, the relation to
% the three interpretations of the expansive study, the connections that are
% the homily's own, the route by which each quoted text reaches the page, and
% the exact loci --- followed by the References.
% Nothing here is read aloud.

\section*{Note on Sources and Delivery}
\runninghead{Sources and Delivery}

{\footnotesize

\textbf{Audience and occasion.} An adult parish assembly, at the Mass
\latin{D.~N. Iesu Christi Regis}, \latin{Dominica ultima octobris}, I classis,
of the 1962 \work{Missale Romanum}, editio typica, pp.~697--698,
nos.~3938--3949, with the proper Preface, p.~291, no.~1080, for the occurrence
of Sunday 25~October 2026, when the feast takes the place of the
Twenty-second Sunday after Pentecost without commemoration. The speech
supposes that it is preached after the Gospel and before the Credo: it speaks
of the Introit, Collect, Epistle and Gradual as heard, and of the Credo,
Offertory, Secret, Preface, Communion and Postcommunion as still to come. The
Alleluia is not used. The votive Tract and the paschal Alleluia printed with
the Mass are not said on the feast and are not used.

\textbf{Length and pace.} The spoken body runs to about 1,530 words, some 11.8
to 12.8 minutes at an unhurried 120 to 130 words a minute by arithmetic alone,
a little beyond the ten to twelve minutes usual for this homily; no one
has delivered or timed it. It was read through on the page for sense and ease
of speech, not aloud.

\textbf{Relation to the study's interpretations.} The homily follows the
second of the full study's three interpretations, \textit{A Kingdom from
Above, Gathered from Every Nation and Served by Kings}: the two wrong answers
to the Gospel's denial, with \work{Quas primas} on the grave error of denying
Christ any authority in civil affairs; Augustine's distinction of ``of'' from
``in'' and of ``from hence'' from ``here'' (his own reading of the Gospel's
words, with the kingdom here until the harvest and sojourning as a stranger),
his kingdom of those who believe, made no longer of the world by rebirth in
Christ, and the Epistle's translation as his account of that change;
Chrysostom's earthly kingship whose strength lies in servants; Aquinas's
kingdom not from here and yet here, because everywhere; the hearers of the truth,
with Augustine's hearing that is believing and obeying, and the grace that
calls; the Collect's nations scattered by the wound of sin and its most loving
rule; the Preface's kingdom of truth; the Secret's gifts of unity and peace;
Augustine's kings bidden not to be sad, for whom it is good to be under
Christ, and his kings adoring among the poor who confess their sins;
Bellarmine's fulfilment of the Gradual's verse, in part in converted princes
and nations and to the letter at the judgment; the Postcommunion's
warfare, to which the homily applies Chrysostom's argument that it is not
Christ's defence and Hilary's kingdom, where sin is overcome and death
destroyed; and Bellarmine's literal sense of the Introit's verse, with the
psalm's justice to the poor. From the first interpretation, \textit{The Slain
Lamb Receives the Kingdom That Was Always His}, it takes Chrysostom's prisoner
made captive by consenting, Aquinas on the Offertory's command to ask, made by
prayer and by passion, Hilary's co-heirs who must become kings, the Lamb
worthy because slain, and the Secret's victim. The difference that matters to
the argument is kept: with Aquinas the nations are the Son's by right before
he asks, and he asks because the Father wills it, so the kingship is not made
to begin in the trial. From the third interpretation, \textit{The Kingdom
Handed Over Is the Redeemed, Brought to the Vision of God}, it takes only
Augustine on the Communion's verse, the Lord dwelling first in his saints in
the flood of this world and afterward reigning in them for ever; the handing
over of the kingdom and its difficulty are left to the studies. Following the
Homiletic Directory, the speech proclaims the kingdom won in Christ's passion
and offered at the altar, leads the assembly toward the Creed and Communion,
and closes its application with one response, said once to the hearers in the
second person and drawn from the Introit's verse as Bellarmine reads it: to
ask God for that judgment and justice in whatever charge they hold, and to
use the charge for the people under it. It ends on its conclusion, that the
King whom Pilate questioned has his throne in those who hear his voice.
These connections are the homily's own and credited to no Father: a kingdom
bounded by the place its beginning comes from, and one from above bounded by
none, drawn from Aquinas's point that the kingdom is everywhere, not in his words; the
Epistle's translation as a carrying out from under the power of darkness of
those who still live in the world; Chrysostom's servants as the reason the
bound King still answers as a king; the truth of the Gospel's last verse as
the first good the Preface names; the trial as the asking of the Offertory,
made by the prisoner's giving himself up; the Secret's unity and peace as the
gathering of the nations the Collect calls scattered, a gift asked for at
this altar; the Postcommunion's war as the war against sin, begun by each in
himself, which neither Chrysostom nor Hilary says of that prayer; the
kingdom's reach to every place where a hearer has charge of others, and the
Introit's verse taken as a prayer for that charge; the Creed as the answer of
those who have heard his voice; and the closing throne in those who hear his
voice, which joins Augustine's kingdom of those who believe to his reading of
the Communion's verse. Augustine's \work{Letter}~185 and Schuster's notes are
not used.

\textbf{Texts quoted, and their route.} Scripture is the Douay--Rheims,
Challoner's revision: Ps 71:8, 11; Ps 2:8; Jn 18:36, 37; Col 1:13; Apoc 5:12;
Ps 71:2; Ps 28:10. The scene is drawn from Jn 18:12, 28--29 and 33--35, and
the psalm's justice to the poor from the rest of Ps 71:2. The Preface's goods
of the kingdom, the Secret's ``this host for the reconciliation of humanity''
and ``will bestow on all Gentiles the gifts of unity and peace'', and the
Postcommunion's ``we who glory to fight under the standard of Christ the
King'' are in the English of the Lasance \work{New Roman Missal} of 1945, a
historical study aid and not an approved liturgical translation of the 1962
Missal; the Collect and the rest of those prayers are described, not quoted.
Augustine and Chrysostom are quoted in the English of the Nicene and
Post-Nicene Fathers, and Bellarmine in O'Sullivan's English. Aquinas and
Hilary are paraphrased from their Latin; Augustine on those who belong to the
kingdom, on hearing the truth, on the kings' good in being under Christ, on
the kings who adore and on the Lord dwelling in his saints, and Bellarmine on
the fulfilments of the Gradual's verse and on David's prayer for Solomon, are paraphrased from
the same English. \work{Quas primas} is described, not quoted, from the Holy
See's English.

\textbf{Exact loci.} Introit, Apoc 5:12; verse Ps 71:1 in the Missal's
numbering (Douay 71:2), no.~3938; Collect, no.~3939; Epistle, Col 1:13,
no.~3940; Gradual, Ps 71:8, 11, no.~3941; Gospel, Jn 18:33--37, no.~3945;
Offertory, Ps 2:8, no.~3946; Secret, no.~3947; Preface, no.~1080; Communion,
Ps 28:10--11 in the Missal's numbering (one Douay verse, 28:10), no.~3948;
Postcommunion, no.~3949. Augustine, \work{Tractates on the Gospel of John}
115.2 (``of'' and ``in'', ``from hence'' and ``here''; the kingdom of those
who believe; Col 1:13), 115.4 (the hearers of the truth); \work{Enarrationes
in Psalmos} 2.9 (the kings of the earth), 28.10 (the Lord dwelling in his
saints), 71.14 (the kings among the poor). Chrysostom, \work{Homilies on
John} 83.4 (the servants who do not fight; the prisoner's consent). Aquinas,
\work{Super Io.}\ 18 lect.~6 (the kingdom not from here and yet here);
\work{Postilla super Psalmos} 2, nn.~5--6 (the asking by prayer and passion).
Hilary, \work{Tractatus super Psalmos} 2.42--43. Bellarmine,
\work{Explanatio}, Ps~71 (on vv.~1 and 11). Pius~XI, \work{Quas primas}
n.~17.

\par}

\Needspace{14\baselineskip}
\section*{References}
\runninghead{References}

{\footnotesize
\begin{itemize}\setlength{\itemsep}{0.15em}\setlength{\parskip}{0pt}
\item \work{Missale Romanum ex decreto Sacrosancti Concilii Tridentini restitutum, Summorum Pontificum cura recognitum}, editio typica (Typis Polyglottis Vaticanis, 1962), CMAA facsimile: Sanctorale, pp.~697--698, nos.~3938--3949; \work{Praefationes}, p.~291, no.~1080.
\item Lasance, F.~X., \work{The New Roman Missal in Latin and English}, revised edition (New York: Benziger Brothers, 1945), pp.~813--814 (Preface of Jesus Christ, King) and 1257--1260 (Feast of Our Lord Jesus Christ, King).
\item \work{The Holy Bible Translated from the Latin Vulgate} (Douay--Rheims, Challoner revision), Project Gutenberg text: Ps 2:8; 28:10; 71:2, 8, 11; Jn 18:12, 28--29, 33--37; Col 1:13; Apoc 5:12.
\item Augustine, \work{Tractates on the Gospel of John} 115.2, 115.4, NPNF, first series, vol.~7 (CCEL text).
\item Augustine, \work{Enarrationes in Psalmos} (\work{Expositions on the Book of Psalms}), NPNF, first series, vol.~8 (CCEL text): 2.9; 28.10; 71.14.
\item Bellarmine, Robert, \work{Explanatio in Psalmos}, in John O'Sullivan's English, \work{A Commentary on the Book of Psalms} (Dublin: James Duffy, 1866), on Ps 71.
\item Chrysostom, John, \work{Homilies on John}, homily 83, NPNF, first series, vol.~14 (New Advent text).
\item Hilary of Poitiers, \work{Tractatus super Psalmos} 2.42--43, ed. A.~Zingerle, CSEL 22 (1891) (Internet Archive page images).
\item Pius~XI, \work{Quas primas}, encyclical of 11~December 1925, English on the Holy See's website, n.~17.
\item Thomas Aquinas, \work{Super Evangelium S.~Ioannis lectura}, chapter 18, lect.~6; \work{Postilla super Psalmos}, on the second psalm, nn.~5--6 (Corpus Thomisticum transcriptions).
\end{itemize}
\par}
